% Supplemental derivation for the finite-index proof in Paper I.
% GPT-6 Astra (OpenAI, ChatGPT Pro mode), effort Pro, 17 November 2026.
\paragraph{The extension transgression, including its sign convention.}
Here is a direct derivation of the transgression used in
\cref{wb:prop:finite-indices}. Write the cyclic quotient additively as
$C_m=\{0,\ldots,m-1\}$ and choose its section $s(a)=g^a$ in
$\Gamma=\langle\Lambda,g\mid g\lambda g^{-1}=A\lambda,
g^m=v\rangle$. Since $Av=v$, its factor set is
\[
 c(a,b)=\left\lfloor\frac{a+b}{m}\right\rfloor v,
 \qquad s(a)s(b)=c(a,b)s(a+b\bmod m).
\]
For an invariant integral covector $\xi$, extend its homomorphism on
$\Lambda$ to the one-cochain $F(\lambda s(a))=\xi(\lambda)$ on
$\Gamma$. The inhomogeneous bar boundary
\cite[Chapter~I, \S5, equation~(5.2)]{Brown1982} gives
\[
 (\delta F)(\lambda s(a),\mu s(b))
 =F(\mu s(b))-F(\lambda s(a)\mu s(b))+F(\lambda s(a))
 =-\xi(c(a,b)).
\]
It is inflated from the displayed cyclic two-cocycle. This is precisely
the obstruction to extending $\xi$ to a cocycle on $\Gamma$, hence the
transgression in the cohomological Hochschild--Serre sequence
\cite[Chapter~VII, \S\S5--6, Theorem~6.3]{Brown1982}.
The carry cocycle $\lfloor(a+b)/m\rfloor$ generates
$H^2(C_m;\mathbb Z)=\mathbb Z/m$: it is the connecting cocycle of
the character $a\mapsto a/m\bmod\mathbb Z$ for
$0\to\mathbb Z\to\mathbb R\to\mathbb R/\mathbb Z\to0$.
Indeed positive-degree cohomology with real coefficients vanishes by
averaging over the finite group, so the connecting map identifies
$\operatorname{Hom}(C_m,\mathbb R/\mathbb Z)$ with
$H^2(C_m;\mathbb Z)$.
If the generator of this last cyclic group is chosen as the negative
of the carry class, the formula is
\[
                 d_2\xi=\xi(v)\bmod m.
\]
This is the convention used in the finite-index proof. Reversing that
generator reverses all these transgression signs together. The image
indices and the nonzero order-two value
$d_2(\gamma\eta_2)=\pm[\eta_2]$ are unchanged. Multiplicativity and
the derivation rule are the filtered cup-product construction of
\cite[Chapter~VII, \S5, pp.~170--171]{Brown1982}.
